\section{Case Studies}
\label{sec:casestudies}

\begin{figure*}[t]
    \centering
    \includegraphics[width=\linewidth]{germany_austria_aliasing_diagnostic.pdf}
    \caption{A majority of foreign identities in Germany and Austria share local identifier suffixes, and the sub-kilometer separation of matched pairs supports identifier aliasing.}
    \label{fig:germany-austria-aliasing}
\end{figure*}

\begin{figure*}[t]
    \centering
    \includegraphics[width=\linewidth]{satellite_world_map.pdf}
    \caption{Satellite-associated cellular observations combine TAC centroids for D2C-like populations with sparse directly assigned and onboard-network locations, while all Starlink and SpaceX observations share one category.}
    \label{fig:satellite-world-map}
\end{figure*}

\begin{figure*}[t]
    \centering
    \setlength{\tabcolsep}{1pt}
    \begin{tabular}{cccc}
        \includegraphics[width=0.245\linewidth]{ocean_vessel_banglalink.pdf} &
        \includegraphics[width=0.245\linewidth]{ocean_vessel_algeria_18413.pdf} &
        \includegraphics[width=0.245\linewidth]{ocean_vessel_algeria_18611.pdf} &
        \includegraphics[width=0.245\linewidth]{ocean_vessel_china_nr_relocation.pdf}
    \end{tabular}
    \caption{Large offshore coordinate changes comprise three temporally plausible vessel-carried GSM candidates and one repeatedly supported China Mobile NR relocation separated by a 129-day observation gap.}
    \label{fig:ocean-vessel-movements}
\end{figure*}

\subsection{Moving Base Stations}

\begin{figure*}[t]
    \centering
    \includegraphics[width=\linewidth]{moving_mccs_overview.pdf}
    \caption{Among cell identities with an exact observed span over 10\,km, most movements are short, while longer-distance endpoints are globally distributed and vary substantially in observation support.}
    \label{fig:moving-mccs-overview}
\end{figure*}
